
\documentclass[12pt,a4paper]{article}

\usepackage[margin=1.4cm]{geometry}
\usepackage{array}
\usepackage{enumitem}
\usepackage{multicol}

% Global compact spacing for lists
\setlist[itemize]{
	topsep=0pt,
	partopsep=0pt,
	parsep=0pt,
	itemsep=0pt
}

\setlist[enumerate]{
	topsep=0pt,
	partopsep=0pt,
	parsep=0pt,
	itemsep=0pt
}

\newcolumntype{L}[1]{>{\raggedright\arraybackslash}p{#1}}

\newcommand{\checkbox}{\fbox{\rule{0pt}{0.8ex}\rule{0.8ex}{0pt}}}

\begin{document}
	
	\begin{center}
		{\Large\bfseries C5: ORGANIZAR LO QUE APRENDIMOS}
		
		\vspace{0.1cm}
		
		Emprendimiento e Innovación Grado 10
		
		\vspace{0.3cm}
		
		\begin{tabular}{L{0.12\textwidth}L{0.25\textwidth}L{0.12\textwidth}L{0.26\textwidth}}
			Grado: & \rule{0.8\linewidth}{0.4pt} & Fecha: & \rule{0.8\linewidth}{0.4pt} \\
			Nombre 1: & \multicolumn{3}{l}{\rule{0.63\textwidth}{0.4pt}} \\
			Nombre 2: & \multicolumn{3}{l}{\rule{0.63\textwidth}{0.4pt}} \\
		\end{tabular}
		
	\end{center}
	\vspace{-7mm}
	\section*{Objetivo}
	Durante las clases anteriores identificaron una situación, observaron lo que sucede en el contexto y escucharon a posibles usuarios. Ahora organizarán la información que recopilaron para construir una representación clara del contexto. \textbf{El objetivo es} encontrar conexiones entre las personas, sus experiencias, el espacio donde ocurre la situación y los hallazgos obtenidos.En esta etapa todavía no necesitan diseñar una solución. Primero deben organizar lo que saben y reconocer qué información es importante.
	
	\vspace{-4mm}
	\section*{Revisen la evidencia}
	Reúnan la información obtenida durante las clases anteriores.Revisen:
	
	\begin{multicols}{2}
		\begin{itemize}[leftmargin=0.6cm,itemsep=0pt]
			\item Las observaciones realizadas.
			\item Las respuestas de los usuarios.
			\item Las necesidades identificadas.
			\item Los comportamientos observados.
			\item Los patrones encontrados.
			\item Las situaciones inesperadas.
			\item Las preguntas que todavía permanecen abiertas.
		\end{itemize}
	\end{multicols}
	
	\textbf{Regla:} Organicen información que provenga de la evidencia recopilada. No agreguen como hechos ideas que todavía no han comprobado.
	
	\vspace{-4mm}
	\section*{1. Seleccionen al usuario principal}
	
	A partir de las observaciones y conversaciones realizadas, identifiquen el grupo de personas que experimenta con mayor frecuencia la situación investigada.
	
	\begin{center}
		\small
		\begin{tabular}{|L{0.55\textwidth}|L{0.4\textwidth}|}
			\hline
			\textbf{Elemento} & \textbf{Nuestra respuesta} \\
			\hline
			¿Qué situación estamos investigando? & \rule{0pt}{0.65cm} \\
			\hline
			¿Quiénes experimentan esta situación? & \rule{0pt}{0.65cm} \\
			\hline
			¿Quiénes la experimentan con mayor frecuencia? & \rule{0pt}{0.65cm} \\
			\hline
			¿Por qué seleccionamos a este usuario o grupo? & \rule{0pt}{0.75cm} \\
			\hline
		\end{tabular}
	\end{center}
	
	\vspace{0.15cm}
	
	\textbf{Importante:} No describan a una persona específica si no es necesario. Pueden representar un usuario típico construido a partir de varias personas que compartan características y experiencias.

	\vspace{-4mm}
	\section*{2. Construyan un perfil de usuario}
	
	Un perfil de usuario resume las características relevantes de las personas que experimentan la situación. No se trata de inventar una historia. El perfil debe representar información que hayan observado o escuchado.
	
	\begin{center}
		\small
		\begin{tabular}{|L{0.22\textwidth}|L{0.2\textwidth}|L{0.51\textwidth}|}
			\hline
			\textbf{Elemento} & \textbf{Nuestro perfil} & \textbf{Ejemplo} \\
			\hline
			Usuario o grupo 
			& \rule{0pt}{0.6cm}
			& Estudiantes que compran alimentos durante el descanso. \\
			\hline
			¿Qué hace habitualmente? 
			& \rule{0pt}{0.8cm}
			& Se dirigen a la cafetería durante el descanso y esperan para comprar. \\
			\hline
			¿Qué necesita realizar? 
			& \rule{0pt}{0.8cm}
			& Comprar alimentos y regresar a sus actividades. \\
			\hline
			¿Qué dificultades experimenta? 
			& \rule{0pt}{0.8cm}
			& En determinados momentos debe esperar durante bastante tiempo. \\
			\hline
			¿Qué le importa en esta situación? 
			& \rule{0pt}{0.8cm}
			& Disponer de suficiente tiempo durante el descanso. \\
			\hline
			¿Qué dijo sobre la situación? 
			& \rule{0pt}{0.8cm}
			& ``A veces paso casi todo el descanso haciendo fila.'' \\
			\hline
		\end{tabular}
	\end{center}
	
	\vspace{-4mm}
	\section*{3. Organicen el contexto}
	
	Ahora representen visualmente dónde ocurre la situación y qué elementos participan en ella. Pueden utilizar un mapa sencillo del espacio, un diagrama de relaciones o una representación del recorrido de las personas. Incluyan:
	
	\begin{multicols}{2}
		\begin{itemize}[leftmargin=0.6cm,itemsep=0pt]
			\item El lugar donde ocurre la situación.
			\item Las personas involucradas.
			\item Los espacios relevantes.
			\item Los objetos o recursos utilizados.
			\item Las acciones principales.
			\item Los puntos donde aparecen dificultades.
			\item Las interacciones importantes.
			\item Las situaciones que se repiten.
		\end{itemize}
	\end{multicols}
	
	\textbf{Nuestro mapa o diagrama del contexto:}
	
	\vspace{0.2cm}
	
	\begin{center}
		\fbox{\rule{0pt}{2cm}\rule{0.84\linewidth}{0pt}}
	\end{center}
	
	\vspace{-4mm}	
	\section*{4. Organicen el recorrido de la situación}
	
	Piensen en lo que ocurre antes, durante y después de la situación investigada.
	
	Identifiquen los momentos en los que las personas realizan acciones, encuentran dificultades o experimentan una necesidad.
	
	\begin{center}
		\small
		\begin{tabular}{|L{0.25\textwidth}|L{0.25\textwidth}|L{0.25\textwidth}|}
			\hline
			\textbf{Antes} & \textbf{Durante} & \textbf{Después} \\
			\hline
			¿Qué sucede? & ¿Qué sucede? & ¿Qué sucede? \\
			& & \\
			\hline
			¿Qué dificultad aparece? & ¿Qué dificultad aparece? & ¿Qué dificultad aparece? \\
			& & \\
			\hline
		\end{tabular}
	\end{center}
	
	\vspace{0.2cm}
	
	\textbf{¿En qué momento parece concentrarse la mayor dificultad y cuál es exactamente?}
	
	\noindent\rule{\linewidth}{0.4pt}
	
	\vspace{0.35cm}
	
	\noindent\rule{\linewidth}{0.4pt}
	
	\section*{5. Organicen los hallazgos}
	
	Revisen toda la información recopilada y seleccionen los hallazgos más importantes. Un hallazgo debe expresar algo que aprendieron a partir de la evidencia.
	
	\textbf{Eviten escribir solamente:} ``Hay un problema con la cafetería.''
	
	\textbf{Es mejor escribir:} ``Observamos que durante los primeros minutos del descanso se forma una fila extensa en la cafetería y varios estudiantes deben esperar antes de poder comprar sus alimentos. Algunos señalaron que la espera reduce el tiempo que tienen disponible para comer y realizar otras actividades durante el descanso.''
	
	\begin{center}
		\small
		\begin{tabular}{|L{0.10\textwidth}|L{0.35\textwidth}|L{0.37\textwidth}|}
			\hline
			\textbf{\#} & \textbf{Hallazgo} & \textbf{Evidencia que lo respalda} \\
			\hline
			1 & & \\
			\hline
			2 & & \\
			\hline
			3 & & \\
			\hline
			4 & & \\
			\hline
			$\vdots$ & $\vdots$ & $\vdots$ \\
			\hline
		\end{tabular}
	\end{center}
	
	\section*{6. Conecten los hallazgos}
	
	Ahora busquen relaciones entre los hallazgos.Pueden agruparlos según:
	\begin{multicols}{2}
		\begin{itemize}[leftmargin=0.6cm,itemsep=0pt]
			\item Personas.
			\item Comportamientos.
			\item Espacios.
			\item Momentos.
			\item Dificultades.
			\item Necesidades.
			\item Factores del contexto.
		\end{itemize}
	\end{multicols}
	
	\textbf{Nuestro listado de categorías:}
	
	\begin{center}
		\small
		\begin{tabular}{|L{0.25\textwidth}|L{0.55\textwidth}|}
			\hline
			\textbf{Categoría} & \textbf{Hallazgos relacionados} \\
			\hline
			Personas & \\
			\hline
			Comportamientos & \\
			\hline
			Espacios / momentos & \\
			\hline
			Dificultades / necesidades & \\
			\hline
		\end{tabular}
	\end{center}
	
	\vspace{0.2cm}
	
	\textbf{¿Qué relación importante encontramos entre nuestros hallazgos?}
	
	\noindent\rule{\linewidth}{0.4pt}
	
	\vspace{0.35cm}
	
	\noindent\rule{\linewidth}{0.4pt}
	
	\section*{7. Identifiquen lo que todavía no saben}
	
	Organizar la información también permite reconocer qué preguntas todavía no pueden responder.
	
	\begin{center}
		\small
		\begin{tabular}{|L{0.40\textwidth}|L{0.4\textwidth}|}
			\hline 
			\textbf{Lo que sabemos} & \textbf{Lo que todavía necesitamos comprender} \\
			\hline
			\rule{0pt}{1.5cm} & \rule{0pt}{1.5cm} \\
			\hline
		\end{tabular}
	\end{center}
	
	\vspace{0.2cm}
	
	\textbf{Pregunta que consideramos más importante investigar:}
	
	\noindent\rule{\linewidth}{0.4pt}
	
	\vspace{0.35cm}
	
	\noindent\rule{\linewidth}{0.4pt}
	
	\section*{8. Construyan nuestro mapa de hallazgos}
	
	Utilicen toda la información anterior para construir una representación final del contexto. El mapa debe mostrar las conexiones entre:
	
	\begin{center}
		\textbf{Usuario}
		\quad $\longleftrightarrow$ \quad
		\textbf{Situación}
		\quad $\longleftrightarrow$ \quad
		\textbf{Contexto}
		\quad $\longleftrightarrow$ \quad
		\textbf{Necesidad}
		\quad $\longleftrightarrow$ \quad
		\textbf{Evidencia}
	\end{center}
	
	No necesitan utilizar un formato específico. Pueden construir un mapa conceptual, diagrama, esquema visual, mapa del contexto o cualquier otra representación que permita comprender la situación.
	
	\textbf{Nuestro mapa de hallazgos:}
	
	\vspace{0.2cm}
	
	\begin{center}
		\fbox{\rule{0pt}{2cm}\rule{0.84\linewidth}{0pt}}
	\end{center}
	
	\section*{Reflexión final}
	
	\textbf{1. ¿Qué aprendimos sobre los usuarios que no era evidente al principio?}
	
	\vspace{0.35cm}
	
	\rule{\linewidth}{0.4pt}
	
	\vspace{0.4cm}
	
	\textbf{2. ¿Qué relación importante encontramos entre las personas, el contexto y la necesidad?}
	
	\vspace{0.35cm}
	
	\rule{\linewidth}{0.4pt}
	
	\vspace{0.4cm}
	
	\textbf{3. ¿Qué hallazgo está respaldado por la evidencia más sólida? ¿Por qué?}
	
	\vspace{0.35cm}
	
	\rule{\linewidth}{0.4pt}
	
	\vspace{0.4cm}
	
	\textbf{4. ¿Qué información todavía necesitamos comprender?}
	
	\vspace{0.35cm}
	
	\rule{\linewidth}{0.4pt}
	
	\vspace{0.4cm}
	
	\textbf{5. Después de organizar la información, ¿cambió nuestra comprensión de la situación? ¿Cómo?}
	
	\vspace{0.35cm}
	
	\rule{\linewidth}{0.4pt}
	
	\section*{Verificación final}
	
	\checkbox\quad Organizamos información obtenida mediante observaciones y conversaciones.
	
	\vspace{0.15cm}
	
	\checkbox\quad Construimos un perfil de usuario basado en evidencia.
	
	\vspace{0.15cm}
	
	\checkbox\quad Representamos el contexto mediante un mapa, esquema o diagrama.
	
	\vspace{0.15cm}
	
	\checkbox\quad Organizamos los principales hallazgos de nuestra investigación.
	
	\vspace{0.15cm}
	
	\checkbox\quad Relacionamos usuarios, comportamientos, espacios y necesidades.
	
	\vspace{0.15cm}
	
	\checkbox\quad Diferenciamos lo que sabemos de lo que todavía necesitamos comprender.
	
	\vspace{0.15cm}
	
	\checkbox\quad Utilizamos evidencia para respaldar nuestros hallazgos.
	
	\vspace{0.15cm}
	
	\checkbox\quad Nuestra representación permite que otra persona comprenda mejor la situación.
	
	\vspace{0.3cm}
	
	\begin{center}
		\textbf{Organicen la evidencia. Encuentren conexiones. Representen lo que aprendieron.}
	\end{center}
	
\end{document}

